\documentclass[11pt]{article}
\usepackage[margin=1in]{geometry}
\usepackage[hidelinks]{hyperref}
\title{Weightlifting: Training Log}
\author{Liam McGrath}
\date{4 October 2026}
\begin{document}
\maketitle

\section{Overview}
Weightlifting since 2012 has been an outlet for getting swole, feeling peak and getting endorphins.

\section{Past peaks}
At the first peak, a 315 lb squat and deadlift could be performed, and at the second peak a 205 lb barbell press. At another peak, 65 lb for 6 reps could be performed on both chest and shoulder press\ldots{} I know, pretty lame. These peaks will be surpassed.

\begin{table}[htbp]
\centering
\begin{tabular}{lll}
\hline
When & Lift & Weight \\
\hline
First peak & Squat & 315 lb \\
First peak & Deadlift & 315 lb \\
Second peak & Barbell press & 205 lb \\
Another peak & Chest press & 65 lb for 6 reps \\
Another peak & Shoulder press & 65 lb for 6 reps \\
\hline
\end{tabular}
\caption{Past peaks, as recalled.}
\end{table}

\section{The current log}
Working sets have been logged since 14 September 2026, with dropsets excluded, on a rotation of chest, back and shoulder days. The Weightlifting page of the website charts the volume of each session and the top set of every exercise logged more than once, with the full log as JSON and CSV downloads.

\end{document}
